% !TEX root = main.tex
% Glossary of Terms (included via \glossaryofterms{% !TEX root = main.tex
% Glossary of Terms (included via \glossaryofterms{% !TEX root = main.tex
% Glossary of Terms (included via \glossaryofterms{% !TEX root = main.tex
% Glossary of Terms (included via \glossaryofterms{\input{listofdefinitions}} in main.tex)
%
% A two-column table: the term, then its definition. List terms in
% alphabetical order. Start each definition with a capital letter and write
% it in sentence case; capitalize only acronyms and proper names (GSO).
% The table continues across pages automatically. To change the width of
% the term column, give a width: \begin{glossarytable}[2in]
%
% Commands cannot be shown with \verb here, because this file is passed as an
% argument; write \texttt{\textbackslash{}cite\{key\}} instead.

\begin{glossarytable}
\term{Argument}{Input given to a command in braces, as in \texttt{\textbackslash{}section\{Title\}}. Optional arguments go in square brackets, as in \texttt{\textbackslash{}includegraphics[width=3in]\{file\}}.}
\term{Auxiliary file}{A file such as \texttt{main.aux} or \texttt{main.toc} that \LaTeX\ writes during one run and reads on the next to resolve cross-references and lists. Deleting these files is a safe first step when a build behaves strangely.}
\term{Biber}{The program that reads \texttt{references.bib} and sorts and formats the bibliography for \texttt{biblatex}. It runs between \LaTeX\ runs.}
\term{\texttt{biblatex}}{The package that formats citations and the bibliography. The citation style is chosen once, in \texttt{packages.tex}.}
\term{\BibTeX\ entry}{One record in a \texttt{.bib} file, such as \texttt{@article\{key, ...\}}. The key is the name used in \texttt{\textbackslash{}cite\{key\}}.}
\term{Class file}{A file ending in \texttt{.cls} that defines a document's overall layout. This template's class, \texttt{njit\_thesis.cls}, implements the NJIT formatting rules.}
\term{Command}{An instruction that begins with a backslash, such as \texttt{\textbackslash{}chapter} or \texttt{\textbackslash{}cite}.}
\term{Comment}{Text after a \texttt{\%} sign on a line. \LaTeX\ ignores it, so comments are useful for notes and for temporarily removing material.}
\term{Cross-reference}{A reference to a numbered item, made with \texttt{\textbackslash{}label} and \texttt{\textbackslash{}ref} (or \texttt{\textbackslash{}eqref} for equations). Numbers update automatically when material moves.}
\term{Environment}{A block of material between \texttt{\textbackslash{}begin\{name\}} and \texttt{\textbackslash{}end\{name\}}, such as \texttt{figure}, \texttt{table}, or \texttt{equation}.}
\term{Float}{A figure or table that \LaTeX\ may move to a nearby page to avoid large gaps. Always refer to floats by number, not as ``the figure below.''}
\term{Label}{A name attached to a numbered item with \texttt{\textbackslash{}label\{name\}}, so it can be referred to with \texttt{\textbackslash{}ref\{name\}}. Prefixes such as \texttt{fig:} or \texttt{eq:} keep labels organized.}
\term{\texttt{latexmk}}{A program that runs \LaTeX\ and Biber as many times as needed to produce a complete document. Running \texttt{latexmk -pdf main.tex} builds this thesis.}
\term{Log file}{The file \texttt{main.log}, which records every warning and error from a run. Check it for undefined references and overfull boxes before submitting.}
\term{Macro}{A command you define yourself with \texttt{\textbackslash{}newcommand}, usually in \texttt{macros.tex}, to keep notation consistent and save typing.}
\term{Overfull box}{A warning that a line or object is too wide and sticks into the margin. The GSO checks margins, so fix these by rewording, adding a hyphenation point, or resizing the object.}
\term{Overleaf}{A website for editing and compiling \LaTeX\ documents in a browser, with no local installation.}
\term{Package}{A file loaded with \texttt{\textbackslash{}usepackage} that adds features, such as \texttt{amsmath} for mathematics or \texttt{graphicx} for images. This template loads its packages in \texttt{packages.tex}.}
\term{Preamble}{Everything in \texttt{main.tex} before \texttt{\textbackslash{}begin\{document\}}, where the class, packages, and macros are loaded.}
\term{Undefined reference}{A \texttt{\textbackslash{}ref} or \texttt{\textbackslash{}cite} that \LaTeX\ cannot resolve, shown as ``??'' in the PDF. It usually means a misspelled label or citation key, or that the document needs another run.}
\end{glossarytable}
} in main.tex)
%
% A two-column table: the term, then its definition. List terms in
% alphabetical order. Start each definition with a capital letter and write
% it in sentence case; capitalize only acronyms and proper names (GSO).
% The table continues across pages automatically. To change the width of
% the term column, give a width: \begin{glossarytable}[2in]
%
% Commands cannot be shown with \verb here, because this file is passed as an
% argument; write \texttt{\textbackslash{}cite\{key\}} instead.

\begin{glossarytable}
\term{Argument}{Input given to a command in braces, as in \texttt{\textbackslash{}section\{Title\}}. Optional arguments go in square brackets, as in \texttt{\textbackslash{}includegraphics[width=3in]\{file\}}.}
\term{Auxiliary file}{A file such as \texttt{main.aux} or \texttt{main.toc} that \LaTeX\ writes during one run and reads on the next to resolve cross-references and lists. Deleting these files is a safe first step when a build behaves strangely.}
\term{Biber}{The program that reads \texttt{references.bib} and sorts and formats the bibliography for \texttt{biblatex}. It runs between \LaTeX\ runs.}
\term{\texttt{biblatex}}{The package that formats citations and the bibliography. The citation style is chosen once, in \texttt{packages.tex}.}
\term{\BibTeX\ entry}{One record in a \texttt{.bib} file, such as \texttt{@article\{key, ...\}}. The key is the name used in \texttt{\textbackslash{}cite\{key\}}.}
\term{Class file}{A file ending in \texttt{.cls} that defines a document's overall layout. This template's class, \texttt{njit\_thesis.cls}, implements the NJIT formatting rules.}
\term{Command}{An instruction that begins with a backslash, such as \texttt{\textbackslash{}chapter} or \texttt{\textbackslash{}cite}.}
\term{Comment}{Text after a \texttt{\%} sign on a line. \LaTeX\ ignores it, so comments are useful for notes and for temporarily removing material.}
\term{Cross-reference}{A reference to a numbered item, made with \texttt{\textbackslash{}label} and \texttt{\textbackslash{}ref} (or \texttt{\textbackslash{}eqref} for equations). Numbers update automatically when material moves.}
\term{Environment}{A block of material between \texttt{\textbackslash{}begin\{name\}} and \texttt{\textbackslash{}end\{name\}}, such as \texttt{figure}, \texttt{table}, or \texttt{equation}.}
\term{Float}{A figure or table that \LaTeX\ may move to a nearby page to avoid large gaps. Always refer to floats by number, not as ``the figure below.''}
\term{Label}{A name attached to a numbered item with \texttt{\textbackslash{}label\{name\}}, so it can be referred to with \texttt{\textbackslash{}ref\{name\}}. Prefixes such as \texttt{fig:} or \texttt{eq:} keep labels organized.}
\term{\texttt{latexmk}}{A program that runs \LaTeX\ and Biber as many times as needed to produce a complete document. Running \texttt{latexmk -pdf main.tex} builds this thesis.}
\term{Log file}{The file \texttt{main.log}, which records every warning and error from a run. Check it for undefined references and overfull boxes before submitting.}
\term{Macro}{A command you define yourself with \texttt{\textbackslash{}newcommand}, usually in \texttt{macros.tex}, to keep notation consistent and save typing.}
\term{Overfull box}{A warning that a line or object is too wide and sticks into the margin. The GSO checks margins, so fix these by rewording, adding a hyphenation point, or resizing the object.}
\term{Overleaf}{A website for editing and compiling \LaTeX\ documents in a browser, with no local installation.}
\term{Package}{A file loaded with \texttt{\textbackslash{}usepackage} that adds features, such as \texttt{amsmath} for mathematics or \texttt{graphicx} for images. This template loads its packages in \texttt{packages.tex}.}
\term{Preamble}{Everything in \texttt{main.tex} before \texttt{\textbackslash{}begin\{document\}}, where the class, packages, and macros are loaded.}
\term{Undefined reference}{A \texttt{\textbackslash{}ref} or \texttt{\textbackslash{}cite} that \LaTeX\ cannot resolve, shown as ``??'' in the PDF. It usually means a misspelled label or citation key, or that the document needs another run.}
\end{glossarytable}
} in main.tex)
%
% A two-column table: the term, then its definition. List terms in
% alphabetical order. Start each definition with a capital letter and write
% it in sentence case; capitalize only acronyms and proper names (GSO).
% The table continues across pages automatically. To change the width of
% the term column, give a width: \begin{glossarytable}[2in]
%
% Commands cannot be shown with \verb here, because this file is passed as an
% argument; write \texttt{\textbackslash{}cite\{key\}} instead.

\begin{glossarytable}
\term{Argument}{Input given to a command in braces, as in \texttt{\textbackslash{}section\{Title\}}. Optional arguments go in square brackets, as in \texttt{\textbackslash{}includegraphics[width=3in]\{file\}}.}
\term{Auxiliary file}{A file such as \texttt{main.aux} or \texttt{main.toc} that \LaTeX\ writes during one run and reads on the next to resolve cross-references and lists. Deleting these files is a safe first step when a build behaves strangely.}
\term{Biber}{The program that reads \texttt{references.bib} and sorts and formats the bibliography for \texttt{biblatex}. It runs between \LaTeX\ runs.}
\term{\texttt{biblatex}}{The package that formats citations and the bibliography. The citation style is chosen once, in \texttt{packages.tex}.}
\term{\BibTeX\ entry}{One record in a \texttt{.bib} file, such as \texttt{@article\{key, ...\}}. The key is the name used in \texttt{\textbackslash{}cite\{key\}}.}
\term{Class file}{A file ending in \texttt{.cls} that defines a document's overall layout. This template's class, \texttt{njit\_thesis.cls}, implements the NJIT formatting rules.}
\term{Command}{An instruction that begins with a backslash, such as \texttt{\textbackslash{}chapter} or \texttt{\textbackslash{}cite}.}
\term{Comment}{Text after a \texttt{\%} sign on a line. \LaTeX\ ignores it, so comments are useful for notes and for temporarily removing material.}
\term{Cross-reference}{A reference to a numbered item, made with \texttt{\textbackslash{}label} and \texttt{\textbackslash{}ref} (or \texttt{\textbackslash{}eqref} for equations). Numbers update automatically when material moves.}
\term{Environment}{A block of material between \texttt{\textbackslash{}begin\{name\}} and \texttt{\textbackslash{}end\{name\}}, such as \texttt{figure}, \texttt{table}, or \texttt{equation}.}
\term{Float}{A figure or table that \LaTeX\ may move to a nearby page to avoid large gaps. Always refer to floats by number, not as ``the figure below.''}
\term{Label}{A name attached to a numbered item with \texttt{\textbackslash{}label\{name\}}, so it can be referred to with \texttt{\textbackslash{}ref\{name\}}. Prefixes such as \texttt{fig:} or \texttt{eq:} keep labels organized.}
\term{\texttt{latexmk}}{A program that runs \LaTeX\ and Biber as many times as needed to produce a complete document. Running \texttt{latexmk -pdf main.tex} builds this thesis.}
\term{Log file}{The file \texttt{main.log}, which records every warning and error from a run. Check it for undefined references and overfull boxes before submitting.}
\term{Macro}{A command you define yourself with \texttt{\textbackslash{}newcommand}, usually in \texttt{macros.tex}, to keep notation consistent and save typing.}
\term{Overfull box}{A warning that a line or object is too wide and sticks into the margin. The GSO checks margins, so fix these by rewording, adding a hyphenation point, or resizing the object.}
\term{Overleaf}{A website for editing and compiling \LaTeX\ documents in a browser, with no local installation.}
\term{Package}{A file loaded with \texttt{\textbackslash{}usepackage} that adds features, such as \texttt{amsmath} for mathematics or \texttt{graphicx} for images. This template loads its packages in \texttt{packages.tex}.}
\term{Preamble}{Everything in \texttt{main.tex} before \texttt{\textbackslash{}begin\{document\}}, where the class, packages, and macros are loaded.}
\term{Undefined reference}{A \texttt{\textbackslash{}ref} or \texttt{\textbackslash{}cite} that \LaTeX\ cannot resolve, shown as ``??'' in the PDF. It usually means a misspelled label or citation key, or that the document needs another run.}
\end{glossarytable}
} in main.tex)
%
% A two-column table: the term, then its definition. List terms in
% alphabetical order. Start each definition with a capital letter and write
% it in sentence case; capitalize only acronyms and proper names (GSO).
% The table continues across pages automatically. To change the width of
% the term column, give a width: \begin{glossarytable}[2in]
%
% Commands cannot be shown with \verb here, because this file is passed as an
% argument; write \texttt{\textbackslash{}cite\{key\}} instead.

\begin{glossarytable}
\term{Argument}{Input given to a command in braces, as in \texttt{\textbackslash{}section\{Title\}}. Optional arguments go in square brackets, as in \texttt{\textbackslash{}includegraphics[width=3in]\{file\}}.}
\term{Auxiliary file}{A file such as \texttt{main.aux} or \texttt{main.toc} that \LaTeX\ writes during one run and reads on the next to resolve cross-references and lists. Deleting these files is a safe first step when a build behaves strangely.}
\term{Biber}{The program that reads \texttt{references.bib} and sorts and formats the bibliography for \texttt{biblatex}. It runs between \LaTeX\ runs.}
\term{\texttt{biblatex}}{The package that formats citations and the bibliography. The citation style is chosen once, in \texttt{packages.tex}.}
\term{\BibTeX\ entry}{One record in a \texttt{.bib} file, such as \texttt{@article\{key, ...\}}. The key is the name used in \texttt{\textbackslash{}cite\{key\}}.}
\term{Class file}{A file ending in \texttt{.cls} that defines a document's overall layout. This template's class, \texttt{njit\_thesis.cls}, implements the NJIT formatting rules.}
\term{Command}{An instruction that begins with a backslash, such as \texttt{\textbackslash{}chapter} or \texttt{\textbackslash{}cite}.}
\term{Comment}{Text after a \texttt{\%} sign on a line. \LaTeX\ ignores it, so comments are useful for notes and for temporarily removing material.}
\term{Cross-reference}{A reference to a numbered item, made with \texttt{\textbackslash{}label} and \texttt{\textbackslash{}ref} (or \texttt{\textbackslash{}eqref} for equations). Numbers update automatically when material moves.}
\term{Environment}{A block of material between \texttt{\textbackslash{}begin\{name\}} and \texttt{\textbackslash{}end\{name\}}, such as \texttt{figure}, \texttt{table}, or \texttt{equation}.}
\term{Float}{A figure or table that \LaTeX\ may move to a nearby page to avoid large gaps. Always refer to floats by number, not as ``the figure below.''}
\term{Label}{A name attached to a numbered item with \texttt{\textbackslash{}label\{name\}}, so it can be referred to with \texttt{\textbackslash{}ref\{name\}}. Prefixes such as \texttt{fig:} or \texttt{eq:} keep labels organized.}
\term{\texttt{latexmk}}{A program that runs \LaTeX\ and Biber as many times as needed to produce a complete document. Running \texttt{latexmk -pdf main.tex} builds this thesis.}
\term{Log file}{The file \texttt{main.log}, which records every warning and error from a run. Check it for undefined references and overfull boxes before submitting.}
\term{Macro}{A command you define yourself with \texttt{\textbackslash{}newcommand}, usually in \texttt{macros.tex}, to keep notation consistent and save typing.}
\term{Overfull box}{A warning that a line or object is too wide and sticks into the margin. The GSO checks margins, so fix these by rewording, adding a hyphenation point, or resizing the object.}
\term{Overleaf}{A website for editing and compiling \LaTeX\ documents in a browser, with no local installation.}
\term{Package}{A file loaded with \texttt{\textbackslash{}usepackage} that adds features, such as \texttt{amsmath} for mathematics or \texttt{graphicx} for images. This template loads its packages in \texttt{packages.tex}.}
\term{Preamble}{Everything in \texttt{main.tex} before \texttt{\textbackslash{}begin\{document\}}, where the class, packages, and macros are loaded.}
\term{Undefined reference}{A \texttt{\textbackslash{}ref} or \texttt{\textbackslash{}cite} that \LaTeX\ cannot resolve, shown as ``??'' in the PDF. It usually means a misspelled label or citation key, or that the document needs another run.}
\end{glossarytable}
